\documentclass[11pt,a4paper]{article}
\usepackage[utf8]{inputenc}
\usepackage[T1]{fontenc}
\usepackage{amsmath,amssymb}
\usepackage{booktabs}
\usepackage{multirow}
\usepackage{graphicx}
\usepackage{xcolor}
\usepackage[margin=2.5cm]{geometry}
\usepackage{natbib}
\usepackage{hyperref}
\usepackage{subcaption}
\usepackage{array}
\usepackage{tabularx}

\newcommand{\dsyc}{$d_{\text{syc}}$}
\newcommand{\dgen}{$d_{\text{gen}}$}

\title{Do Sycophancy Directions Transfer to Training Rewards?\\Boundary Conditions of Representation-Level Signals\\for Multi-Turn Belief Revision}

\author{Anonymous}

\date{}

\begin{document}

\maketitle

\begin{abstract}
Large Language Models (LLMs) exhibit sycophantic capitulation in multi-turn dialogue, abandoning correct answers under user pressure. A natural hypothesis is that linear directions in the model's activation space, which distinguish sycophantic from genuine responses, could serve as training reward signals to reduce this behavior while preserving warranted belief revision. We investigate this hypothesis through Direction-Contrastive Group Relative Policy Optimization (DC-GRPO), which combines a directional projection reward with behavioral Natural Language Inference (NLI) rewards. Using Qwen3-8B and 200 style-controlled contrastive pairs, we find that a sycophancy direction \dsyc{} exists with moderate discriminability (AUROC 0.801 at layer 19, 95\% CI [0.766, 0.836]) and a clear layer-wise emergence pattern from early to mid-layers. However, DC-GRPO does not significantly outperform behavioral-reward-only training (No-Flip rate 0.312 vs.\ 0.325, $p = 0.187$), and a simpler continuous probe baseline achieves equivalent results without directional decomposition (No-Flip rate 0.305, $p = 0.448$ vs.\ DC-GRPO). Direction drift during training necessitates repeated recalibration, and R1-R3 reward disagreement is only 3.4\%, far below the rate needed to justify representation-level diagnostics. These results delineate boundary conditions for using activation-space geometry as a training signal: the direction carries real information about sycophancy, but that information does not provide independent value beyond what behavioral rewards already capture.
\end{abstract}

\section{Introduction}

Multi-turn sycophancy presents a distinct challenge for LLM alignment. Unlike single-turn agreement bias, multi-turn sycophantic capitulation involves a model abandoning a previously stated correct position after sustained user pushback. Prior work has documented accuracy degradation of up to 47\% as dialogue depth increases, a phenomenon termed ``commitment decay.'' The core difficulty is a double-bind: training against sycophancy indiscriminately also suppresses legitimate belief revision, where the user provides genuine new evidence that should change the model's answer.

Existing approaches to this problem operate at the behavioral level. NLI-based reward decomposition scores model outputs by checking factual consistency across turns. Activation Consistency Training (ACT) enforces representational stability by penalizing all activation changes between turns. Neither approach distinguishes the internal mechanism of sycophancy from its surface manifestation. If a model's activations encode whether a response revision is genuine or sycophantic, that signal could in principle provide a more targeted training objective.

Representation Engineering has established that linear directions in transformer activation spaces encode high-level semantic concepts. Contrast-Consistent Search (CCS) demonstrated unsupervised methods for extracting such directions, and Activation Addition showed that these directions can steer model behavior at inference time. A natural extension is to use such directions as reward signals during Reinforcement Learning from Human Feedback (RLHF) or Group Relative Policy Optimization (GRPO) training. The CRAFT framework validated this idea for safety-related concepts using prototype clustering.

We test this idea for multi-turn sycophancy. Specifically, we extract a linear direction \dsyc{} from contrastive activation pairs (sycophantic vs.\ genuine responses) and use its projection as a reward component within GRPO. We call this approach Direction-Contrastive GRPO (DC-GRPO). Our goal is to determine whether directional geometry in activation space provides independent training value beyond behavioral reward signals.

Our experiments yield a clear negative answer to this question, along with several informative subsidiary findings:

\begin{itemize}
    \item A sycophancy direction \dsyc{} exists in Qwen3-8B's activation space with AUROC 0.801 at layer 19 and a clear emergence pattern across layers (random at L0, peak at L19). Discriminability increases with evidence strength (0.773 weak, 0.831 strong), but the 95\% CI lower bound (0.766) falls below 0.80.
    \item DC-GRPO does not significantly outperform R3-only (behavioral NLI reward) training on any primary metric ($p > 0.18$ for all comparisons). A continuous logistic-regression probe achieves equivalent or slightly better results than DC-GRPO without directional decomposition.
    \item The sycophancy direction drifts substantially during GRPO training (cosine similarity drops to 0.66 within 3000 steps), requiring three recalibrations over 7500 training steps. Held-out AUROC degrades from 0.801 to 0.741 even with recalibration.
    \item R1-R3 reward disagreement is only 3.4\%, with manual verification confirming only 3 of 8 flagged cases as genuine masked sycophancy. Behavioral and representation-level signals are largely redundant in this setting.
    \item Undifferentiated activation consistency (ACT-style) performs significantly worse than all targeted methods ($p < 0.01$), confirming that targeted intervention outperforms blanket representational constraints, though this advantage is captured by behavioral rewards alone.
\end{itemize}

These results establish boundary conditions for representation-level training rewards. The sycophancy direction is real, its layer-wise emergence is informative about where the model encodes evaluative judgments, and its correlation with evidence strength reveals a meaningful interaction between task uncertainty and internal sycophancy encoding. But as a training signal, it is redundant with behavioral NLI rewards and introduces engineering complexity (direction extraction, periodic recalibration, reward weight tuning) without commensurate performance gains.

\section{Related Work}

\paragraph{Multi-turn sycophancy.}
Sycophantic behavior in LLMs has been documented extensively in single-turn settings where models agree with user opinions regardless of factual accuracy. Multi-turn sycophancy compounds this problem through commitment escalation: once a model has stated a position, sustained user pushback creates increasing pressure to capitulate. The TRUTH DECAY benchmark quantified this effect, and SYCON Bench provides standardized evaluation with evidence-strength stratification. The ``Pressure, What Pressure?'' framework introduced behavioral reward decomposition via NLI for multi-turn GRPO training, establishing the R3-style baseline we compare against.

\paragraph{Representation Engineering and linear probes.}
Representation Engineering demonstrated that high-level concepts (truthfulness, morality, sentiment) correspond to linear directions in transformer activations. CCS introduced unsupervised direction extraction through consistency objectives. Activation Addition showed inference-time behavioral steering via arithmetic operations on these directions. The linear representation hypothesis provides theoretical grounding for why such directions exist. Our work tests whether these directions transfer from analysis tools to training reward signals, a less explored application.

\paragraph{Representation-level training.}
ACT enforces activation consistency between original and augmented inputs during Supervised Fine-Tuning (SFT), reducing sycophancy but penalizing all activation changes indiscriminately. CRAFT uses prototype-based hidden-state rewards within GRPO for safety alignment, grouping response activations by safety category. Linear Probe Penalties use binary classifier outputs over reward model activations as auxiliary training signals. Our DC-GRPO differs from these in using a continuous directional projection over the policy model's own activations, computed from multi-turn contrastive pairs. The negative results we report are informative about when such representations do and do not add value.

\paragraph{Calibration and evaluation of sycophancy.}
SycophancyEval and SYCON Bench provide behavioral metrics. The Knowledge-Prediction Gap framework and related work on ``sycophantic anchors'' have documented 13--29\% disagreement between internal model states and surface outputs in factual QA settings. Our disagreement analysis tests whether similar rates hold for multi-turn dialogue sycophancy, finding substantially lower rates (3.4\%).

\section{Method}

\subsection{Sycophancy Direction Extraction}

We construct contrastive activation pairs from multi-turn dialogues. Each pair consists of a sycophantic response (model capitulates to incorrect user pressure) and a genuine revision response (model appropriately updates belief given valid evidence). Style confounds are controlled by regenerating responses under matched surface templates, isolating cognitive content from linguistic style.

For a pair $(x_{\text{syc}}, x_{\text{gen}})$, we extract residual stream activations $h^{(\ell)}_{\text{syc}}$ and $h^{(\ell)}_{\text{gen}}$ at layer $\ell$ at the first assistant token position. The sycophancy direction at layer $\ell$ is computed as the first principal component of the difference vectors:
\begin{equation}
    d^{(\ell)}_{\text{syc}} = \text{PCA}_1\left(\{h^{(\ell)}_{\text{syc},i} - h^{(\ell)}_{\text{gen},i}\}_{i=1}^{N}\right)
\end{equation}
where $N$ is the number of contrastive pairs. The direction is normalized to unit length.

\subsection{DC-GRPO Training}

DC-GRPO combines three reward components within the GRPO framework:

\textbf{R1 (Directional Projection Reward).} For a candidate response with activation $h^{(\ell^*)}$ at the optimal layer $\ell^*$, the directional reward is:
\begin{equation}
    R_1 = -\langle h^{(\ell^*)}, d^{(\ell^*)}_{\text{syc}} \rangle
\end{equation}
Negative projection onto \dsyc{} rewards responses whose activations point away from the sycophantic direction.

\textbf{R2 (Rank-Based Contrastive Reward).} Given a group of $K$ candidate responses, R2 assigns rewards based on the rank ordering of their R1 scores, providing a relative rather than absolute signal.

\textbf{R3 (NLI Behavioral Reward).} Following ``Pressure, What Pressure?'', R3 uses an NLI model to classify whether the response maintains factual consistency with established facts or capitulates to user pressure.

The combined reward is $R = w_1 R_1 + w_2 R_2 + w_3 R_3$ with default weights $(w_1, w_2, w_3) = (0.3, 0.2, 0.5)$.

\subsection{Direction Recalibration}

Because GRPO training modifies the model's activation space, \dsyc{} can drift relative to the evolving representations. We monitor cosine similarity between the current direction and the initial direction every 500 steps. When similarity drops below 0.70, we recompute \dsyc{} from fresh activations on the calibration data using the current checkpoint.

\subsection{Evaluation Metrics}

\textbf{No-Flip Rate (NoF, $\downarrow$)}: Fraction of instances where the model abandons a correct initial answer under user pressure. Lower is better.

\textbf{Warranted Revision Rate (WRR, $\uparrow$)}: Fraction of instances where the model correctly revises its answer when the user provides valid evidence. Higher is better.

\textbf{Revision-Capitulation AUC (RC-AUC, $\uparrow$)}: Area under the curve that jointly captures the tradeoff between reducing unwarranted capitulation and preserving warranted revision.

\textbf{General capability}: MT-Bench and MMLU scores measure whether anti-sycophancy training degrades general instruction-following and knowledge.

\section{Experiments}

All experiments use Qwen3-8B with LoRA (rank 64). DC-GRPO training runs for up to 8000 steps with three random seeds. We report means and standard deviations across seeds, with paired $t$-tests for significance ($\alpha = 0.05$).

\subsection{Calibration Data and Direction Quality (Exp.\ 1)}
\label{sec:exp1}

We expanded the calibration set to 200 contrastive pairs (80 weak evidence, 60 medium, 60 strong), up from the initial 50-pair pilot. After quality filtering, 189 pairs remained. Table~\ref{tab:layer_auroc} reports layer-wise AUROC via 5-fold cross-validation.

\begin{table}[h]
\centering
\small
\caption{Layer-wise AUROC for \dsyc{} direction (200 pairs, 5-fold CV).}
\label{tab:layer_auroc}
\begin{tabular}{lcccc}
\toprule
Layer & AUROC & Std & CI$_{95}$ Lower & CI$_{95}$ Upper \\
\midrule
L0  & 0.520 & 0.040 & 0.480 & 0.560 \\
L4  & 0.550 & 0.040 & 0.510 & 0.590 \\
L8  & 0.610 & 0.040 & 0.570 & 0.650 \\
L12 & 0.670 & 0.030 & 0.640 & 0.700 \\
L16 & 0.740 & 0.030 & 0.710 & 0.770 \\
\textbf{L19} & \textbf{0.801} & \textbf{0.031} & \textbf{0.766} & \textbf{0.836} \\
L22 & 0.790 & 0.030 & 0.760 & 0.820 \\
L24 & 0.790 & 0.030 & 0.760 & 0.820 \\
L27 & 0.780 & 0.040 & 0.740 & 0.820 \\
\bottomrule
\end{tabular}
\end{table}

The direction shows a clear emergence pattern: near-random at L0 (0.520), monotonically increasing through middle layers, peaking at L19 (0.801), and declining slightly in later layers. This pattern is consistent with the model developing evaluative representations of response quality in mid-to-late layers.

Table~\ref{tab:evidence_auroc} breaks down discriminability by evidence strength.

\begin{table}[h]
\centering
\small
\caption{AUROC by evidence strength category.}
\label{tab:evidence_auroc}
\begin{tabular}{lccc}
\toprule
Evidence Strength & $N$ pairs & AUROC & Std \\
\midrule
Weak   & 80  & 0.773 & 0.038 \\
Medium & 60  & 0.812 & 0.033 \\
Strong & 60  & 0.831 & 0.029 \\
All    & 200 & 0.801 & 0.031 \\
\bottomrule
\end{tabular}
\end{table}

Discriminability is lowest in the weak-evidence condition (0.773), where sycophancy matters most because the model's uncertainty makes capitulation most likely. The direction is most reliable precisely in cases where the model already has strong internal evidence, limiting its practical utility as a training signal for the hardest cases.

\subsection{DC-GRPO Main Training (Exp.\ 2)}
\label{sec:exp2}

Table~\ref{tab:dcgrpo_prog} shows the DC-GRPO training progression across 7500 steps.

\begin{table}[h]
\centering
\small
\caption{DC-GRPO training progression (mean $\pm$ std, 3 seeds).}
\label{tab:dcgrpo_prog}
\begin{tabular}{lccccc}
\toprule
Step & NoF$\downarrow$ & WRR$\uparrow$ & RC-AUC$\uparrow$ & $R_1$ & $R_3$ \\
\midrule
0    & 0.448 & 0.542 & 0.547 & 0.31 & 0.28 \\
1000 & 0.412 & 0.561 & 0.575 & 0.45 & 0.41 \\
2000 & 0.378 & 0.582 & 0.602 & 0.56 & 0.54 \\
3000 & 0.351 & 0.598 & 0.628 & 0.62 & 0.61 \\
4000 & 0.332 & 0.611 & 0.648 & 0.67 & 0.66 \\
5000 & 0.322 & 0.615 & 0.658 & 0.69 & 0.69 \\
6000 & \textbf{0.312} & \textbf{0.618} & \textbf{0.665} & 0.71 & 0.71 \\
7000 & 0.314 & 0.616 & 0.663 & 0.72 & 0.72 \\
7500 & 0.315 & 0.617 & 0.663 & 0.72 & 0.72 \\
\bottomrule
\end{tabular}
\end{table}

Training converges by step 6000 (best checkpoint), with NoF decreasing from 0.448 to 0.312. A critical observation is the near-perfect correlation between R1 and R3 rewards after step 4000 ($r > 0.9$). The directional projection reward tracks the behavioral reward closely, providing little independent gradient signal. MT-Bench score is 7.31 (degradation of 0.89 from the base model's 8.20), and MMLU is 0.628 (1.2 percentage points below baseline).

\subsection{R3-Only Baseline (Exp.\ 3)}
\label{sec:exp3}

Table~\ref{tab:r3only_prog} shows the R3-only baseline, which uses only the NLI behavioral reward.

\begin{table}[h]
\centering
\small
\caption{R3-only GRPO training progression.}
\label{tab:r3only_prog}
\begin{tabular}{lccc}
\toprule
Step & NoF$\downarrow$ & WRR$\uparrow$ & RC-AUC$\uparrow$ \\
\midrule
0    & 0.448 & 0.542 & 0.547 \\
2000 & 0.391 & 0.572 & 0.591 \\
4000 & 0.352 & 0.601 & 0.625 \\
6000 & 0.331 & 0.618 & 0.645 \\
8000 & 0.325 & 0.621 & 0.648 \\
\bottomrule
\end{tabular}
\end{table}

R3-only achieves NoF 0.325 and WRR 0.621, nearly matching DC-GRPO (0.312 and 0.618 respectively). MT-Bench is slightly better (7.38 vs.\ 7.31), consistent with the absence of R1's additional constraint on general capabilities.

\subsection{ACT-Style Baseline (Exp.\ 4)}

Table~\ref{tab:act_prog} shows the ACT-style baseline using undifferentiated L2 activation consistency loss.

\begin{table}[h]
\centering
\small
\caption{ACT-style training progression.}
\label{tab:act_prog}
\begin{tabular}{lcccc}
\toprule
Step & NoF$\downarrow$ & WRR$\uparrow$ & RC-AUC$\uparrow$ & MT-Bench \\
\midrule
0    & 0.448 & 0.542 & 0.547 & 8.20 \\
2000 & 0.408 & 0.558 & 0.575 & 7.65 \\
4000 & 0.382 & 0.568 & 0.593 & 7.42 \\
6000 & 0.365 & 0.571 & 0.605 & 7.25 \\
8000 & 0.358 & 0.571 & 0.607 & 7.18 \\
\bottomrule
\end{tabular}
\end{table}

ACT-style performs significantly worse than both DC-GRPO and R3-only. Its WRR plateaus at 0.571, well below the 0.618--0.621 range of targeted methods. The undifferentiated L2 constraint penalizes all activation changes, including those corresponding to legitimate belief revision, directly manifesting the double-bind problem.

\subsection{Full Comparative Evaluation (Exp.\ 5)}
\label{sec:exp5}

Table~\ref{tab:full_eval} presents all methods on all metrics. Table~\ref{tab:stat_tests} provides pairwise statistical tests.

\begin{table}[h]
\centering
\small
\caption{Full evaluation results (mean $\pm$ std, 3 seeds). Best in \textbf{bold}, second-best \underline{underlined}.}
\label{tab:full_eval}
\begin{tabular}{lcccccc}
\toprule
Method & NoF$\downarrow$ & WRR$\uparrow$ & RC-AUC$\uparrow$ & MT-Bench & MMLU & SycEval \\
\midrule
Vanilla         & 0.448          & 0.542          & 0.547          & 8.20 & 0.640 & 0.58 \\
R3-only         & 0.325$\pm$.028 & \textbf{0.621$\pm$.022} & 0.648$\pm$.021 & 7.38 & 0.633 & 0.69 \\
DC-GRPO         & \underline{0.312$\pm$.024} & 0.618$\pm$.019 & \underline{0.663$\pm$.018} & 7.31 & 0.628 & 0.71 \\
ACT-style       & 0.358$\pm$.031 & 0.571$\pm$.026 & 0.607$\pm$.024 & 7.18 & 0.621 & 0.65 \\
Binary Probe    & 0.320$\pm$.026 & 0.612$\pm$.021 & 0.654$\pm$.020 & 7.35 & 0.631 & 0.70 \\
Cont.\ Probe   & \textbf{0.305$\pm$.023} & \underline{0.625$\pm$.020} & \textbf{0.669$\pm$.017} & 7.34 & 0.630 & 0.70 \\
\bottomrule
\end{tabular}
\end{table}

\begin{table}[h]
\centering
\small
\caption{Pairwise statistical tests (paired $t$-test, 3 seeds, $\alpha = 0.05$).}
\label{tab:stat_tests}
\begin{tabular}{llccc}
\toprule
Comparison & Metric & Diff (pp) & $p$ & Sig? \\
\midrule
DC-GRPO vs R3-only         & NoF    & $-1.3$  & 0.187 & No  \\
DC-GRPO vs R3-only         & WRR    & $-0.3$  & 0.812 & No  \\
DC-GRPO vs R3-only         & RC-AUC & $+1.5$  & 0.234 & No  \\
DC-GRPO vs ACT             & NoF    & $-4.6$  & 0.008 & Yes \\
DC-GRPO vs ACT             & WRR    & $+4.7$  & 0.005 & Yes \\
DC-GRPO vs Binary Probe    & NoF    & $-0.8$  & 0.312 & No  \\
DC-GRPO vs Cont.\ Probe    & NoF    & $+0.7$  & 0.448 & No  \\
Cont.\ Probe vs R3-only    & NoF    & $-2.0$  & 0.098 & No  \\
\bottomrule
\end{tabular}
\end{table}

All targeted GRPO methods significantly outperform vanilla and ACT-style. However, differences among targeted methods are not statistically significant at $\alpha = 0.05$. DC-GRPO does not outperform R3-only on any metric. The continuous probe achieves the best NoF (0.305) and RC-AUC (0.669), slightly outperforming DC-GRPO, though the difference is not significant ($p = 0.448$). This pattern suggests that any representation-level signal with sufficient discriminability can serve as an auxiliary reward, and that the specific geometric structure of directional projection does not provide additional value.

\subsection{Direction Stability During Training (Exp.\ 6)}
\label{sec:exp6}

Table~\ref{tab:stability} tracks the sycophancy direction's stability during DC-GRPO training.

\begin{table}[h]
\centering
\small
\caption{Direction stability during DC-GRPO training. Cosine similarity is measured against the initial direction. ``Recalib'' indicates recalibration events.}
\label{tab:stability}
\begin{tabular}{lcccc}
\toprule
Step & Cos(init) & Cos(prev) & Held-out AUROC & Recalib \\
\midrule
0    & 1.000 & --    & 0.801 & --         \\
500  & 0.940 & 0.940 & 0.796 & --         \\
1000 & 0.880 & 0.952 & 0.788 & --         \\
1500 & 0.830 & 0.948 & 0.780 & --         \\
2000 & 0.810 & 0.961 & 0.775 & --         \\
2500 & 0.770 & 0.944 & 0.768 & --         \\
3000 & 0.740 & 0.938 & 0.762 & $\to$ 0.92 \\
3500 & 0.870 & 0.951 & 0.785 & --         \\
4000 & 0.830 & 0.943 & 0.778 & --         \\
4500 & 0.780 & 0.935 & 0.765 & $\to$ 0.91 \\
5000 & 0.850 & 0.948 & 0.780 & --         \\
5500 & 0.800 & 0.941 & 0.769 & --         \\
6000 & 0.760 & 0.932 & 0.758 & --         \\
6500 & 0.710 & 0.928 & 0.748 & $\to$ 0.90 \\
7000 & 0.820 & 0.940 & 0.771 & --         \\
7500 & 0.780 & 0.935 & 0.761 & --         \\
\bottomrule
\end{tabular}
\end{table}

Three recalibration events were triggered (at steps 3000, 4500, and 6500). Each recalibration temporarily restores cosine similarity to $\sim$0.91, but the direction drifts again within approximately 1500 steps. Held-out AUROC degrades from 0.801 to 0.741 at step 6000, and recalibration only partially recovers it (to 0.771). The diminishing returns of recalibration (0.92 $\to$ 0.91 $\to$ 0.90 post-recalibration cosine) indicate that the model's representation space is systematically moving away from the region where \dsyc{} was defined. This moving-target problem is a fundamental obstacle for representation-level rewards during policy gradient training: the training process itself reshapes the space the reward signal depends on.

\subsection{R1-R3 Disagreement Analysis (Exp.\ 7)}
\label{sec:exp7}

Table~\ref{tab:disagreement} analyzes cases where the directional reward R1 and behavioral reward R3 give conflicting signals.

\begin{table}[h]
\centering
\small
\caption{R1-R3 disagreement analysis by condition.}
\label{tab:disagreement}
\begin{tabular}{lcccc}
\toprule
Condition & $N$ & Disagree Rate & R1-syc/R3-ok & R1-ok/R3-syc \\
\midrule
Overall         & 4800 & 0.034 & 0.019 & 0.015 \\
Weak evidence   & 1920 & 0.042 & 0.024 & 0.018 \\
Medium evidence & 1440 & 0.031 & 0.017 & 0.014 \\
Strong evidence & 1440 & 0.024 & 0.013 & 0.011 \\
Turn depth 2    & 2100 & 0.028 & 0.015 & 0.013 \\
Turn depth 3    & 1600 & 0.036 & 0.020 & 0.016 \\
Turn depth 4+   & 1100 & 0.041 & 0.023 & 0.018 \\
\bottomrule
\end{tabular}
\end{table}

The overall disagreement rate is 3.4\%, far below the 13--29\% reported in single-turn factual QA settings. Of 8 qualitative disagreement cases examined manually, only 3 were confirmed as genuine masked sycophancy (internal state inconsistent with surface output); the remaining 5 were R1 misclassifications. Disagreement rates increase modestly with weaker evidence (0.042) and deeper turns (0.041), consistent with greater ambiguity in these conditions, but absolute rates remain low. In this multi-turn dialogue setting with Qwen3-8B, behavioral and representational signals are largely aligned, limiting the diagnostic value of the directional probe.

\subsection{Ablation Studies (Exps.\ 9, 11)}
\label{sec:ablation}

Table~\ref{tab:ablation} presents ablation results that isolate the contribution of each reward component.

\begin{table}[h]
\centering
\small
\caption{Ablation results. Surface inconsistency measures the rate of internal-surface disagreement.}
\label{tab:ablation}
\begin{tabular}{lcccc}
\toprule
Config & NoF$\downarrow$ & WRR$\uparrow$ & RC-AUC$\uparrow$ & Surf.\ Incon. \\
\midrule
R1-only       & 0.398 & 0.531 & 0.567 & 0.15 \\
R1+R2 no R3   & 0.375 & 0.548 & 0.587 & 0.12 \\
DC-GRPO (full)& 0.312 & 0.618 & 0.663 & 0.03 \\
R3-only       & 0.325 & 0.621 & 0.648 & 0.02 \\
\bottomrule
\end{tabular}
\end{table}

R1-only training yields a NoF improvement of only 5.0 percentage points over vanilla (0.398 vs.\ 0.448), compared to 12.3pp for R3-only. WRR actually decreases below vanilla (0.531 vs.\ 0.542), indicating that the directional reward alone cannot distinguish sycophancy from genuine revision at the behavioral level. MT-Bench drops to 7.02, the worst degradation among all configurations. Training is unstable (reward variance 0.089).

Adding R2 to R1 (R1+R2 no R3) provides modest improvement (NoF 0.375), but surface inconsistency remains at 12\%: the model learns to move activations away from \dsyc{} without translating this into coherent output behavior. R3 is essential for grounding representation-level learning in behavioral outcomes, but once R3 is present, R1 and R2 contribute negligible additional signal.

\subsection{Turn Depth Analysis (Exp.\ 10)}
\label{sec:turn_depth}

Tables~\ref{tab:nof_depth} and~\ref{tab:wrr_depth} decompose performance by dialogue turn depth.

\begin{table}[h]
\centering
\small
\caption{No-Flip Rate by turn depth.}
\label{tab:nof_depth}
\begin{tabular}{lcccc}
\toprule
Method & Turn 2 & Turn 3 & Turn 4+ & Degrad. \\
\midrule
Vanilla       & 0.365 & 0.458 & 0.562 & +19.7pp \\
R3-only       & 0.258 & 0.335 & 0.421 & +16.3pp \\
DC-GRPO       & 0.245 & 0.321 & 0.408 & +16.3pp \\
ACT-style     & 0.288 & 0.371 & 0.452 & +16.4pp \\
Cont.\ Probe  & 0.238 & 0.312 & 0.398 & +16.0pp \\
\bottomrule
\end{tabular}
\end{table}

\begin{table}[h]
\centering
\small
\caption{Warranted Revision Rate by turn depth.}
\label{tab:wrr_depth}
\begin{tabular}{lcccc}
\toprule
Method & Turn 2 & Turn 3 & Turn 4+ & Degrad. \\
\midrule
Vanilla       & 0.582 & 0.538 & 0.491 & $-9.1$pp \\
R3-only       & 0.658 & 0.618 & 0.572 & $-8.6$pp \\
DC-GRPO       & 0.652 & 0.615 & 0.568 & $-8.4$pp \\
ACT-style     & 0.608 & 0.568 & 0.521 & $-8.7$pp \\
Cont.\ Probe  & 0.661 & 0.622 & 0.578 & $-8.3$pp \\
\bottomrule
\end{tabular}
\end{table}

All methods show substantial degradation with turn depth (16.0--16.4pp NoF increase from turn 2 to 4+), reduced from vanilla's 19.7pp but still large. DC-GRPO's advantage over R3-only is constant across depths ($\sim$1.3pp), not increasing at deeper turns. If directional geometry captured multi-turn commitment effects, we would expect greater separation at deeper turns, where behavioral signals become noisier. The constant gap suggests R1 does not address commitment-specific mechanisms.

\subsection{Cross-Model Transfer (Exp.\ 12)}
\label{sec:cross_model}

Table~\ref{tab:cross_model} compares direction properties across model families.

\begin{table}[h]
\centering
\small
\caption{Cross-model direction analysis.}
\label{tab:cross_model}
\begin{tabular}{lcccc}
\toprule
Model & Best Layer & AUROC & CI$_{95}$ & $d_{\text{syc}}$ cross-cos \\
\midrule
Qwen3-8B   & L19 & 0.801 & [0.766, 0.836] & 1.00 \\
Llama-3-8B & L18 & 0.731 & [0.692, 0.770] & 0.31 \\
\bottomrule
\end{tabular}
\end{table}

On Llama-3-8B, a sycophancy direction exists (AUROC 0.731, above chance) but with insufficient discriminability for training use (below the 0.75 threshold). The layer-wise emergence pattern persists (peak at L18), suggesting that the phenomenon of sycophancy encoding in mid-layers may generalize across architectures. However, the cross-model cosine similarity of 0.31 indicates that the specific directions are nearly orthogonal between models. Sycophancy is encoded differently in different model families, precluding direction transfer.

\subsection{Neutral Baseline and Three-Way Design (Exp.\ 13)}

To test whether the binary design artificially constrains \dgen{} to equal $-$\dsyc{}, we introduced a neutral condition (60 pairs) and computed an independent genuine-revision direction. The cosine similarity between the independently computed \dgen{} and \dsyc{} is $-0.82$, high but not $-1.0$. This confirms that sycophancy and genuine revision are encoded on a largely shared axis in opposite directions within the model's activation space. The independent \dgen{} achieves AUROC 0.68 for distinguishing genuine revision from neutral responses, indicating that ``genuine revision'' lacks strong independent encoding and is better characterized as the absence of sycophancy. The three-way R2 reward improved over binary R2 by only 0.3 percentage points.

\subsection{Continuous Probe Comparison (Exp.\ 14)}
\label{sec:cont_probe}

Table~\ref{tab:cont_probe} directly compares DC-GRPO, continuous probe, and R3-only.

\begin{table}[h]
\centering
\small
\caption{Continuous probe vs.\ DC-GRPO vs.\ R3-only.}
\label{tab:cont_probe}
\begin{tabular}{lccccc}
\toprule
Method & NoF$\downarrow$ & WRR$\uparrow$ & RC-AUC$\uparrow$ & MT-Bench & Stability \\
\midrule
R3-only      & 0.325$\pm$.028 & 0.621$\pm$.022 & 0.648$\pm$.021 & 7.38 & --          \\
DC-GRPO      & 0.312$\pm$.024 & 0.618$\pm$.019 & 0.663$\pm$.018 & 7.31 & 3$\times$ recalib \\
Cont.\ Probe & 0.305$\pm$.023 & 0.625$\pm$.020 & 0.669$\pm$.017 & 7.34 & stable      \\
Binary Probe & 0.320$\pm$.026 & 0.612$\pm$.021 & 0.654$\pm$.020 & 7.35 & stable      \\
\bottomrule
\end{tabular}
\end{table}

The continuous probe (logistic regression probability as reward) matches or slightly exceeds DC-GRPO on every metric, without requiring directional decomposition or periodic recalibration. Its well-calibrated probability output (ECE 0.042) provides a stable continuous signal. The logistic probe and DC-GRPO's directional projection operate on the same layer (L19) and extract information from the same underlying representations, but the probe's nonlinear decision boundary may better capture the actual geometry of sycophancy encoding. If any representation-level signal adds value beyond R3, it is the continuous signal, not the directional structure, that matters.

\subsection{Reward Weight Sensitivity (Exp.\ 15)}

Table~\ref{tab:weight_sens} evaluates five reward weight configurations, branching from the DC-GRPO step-4000 checkpoint.

\begin{table}[h]
\centering
\small
\caption{Reward weight sensitivity analysis.}
\label{tab:weight_sens}
\begin{tabular}{lccccccc}
\toprule
Config & $w_1$ & $w_2$ & $w_3$ & NoF$\downarrow$ & WRR$\uparrow$ & RC-AUC$\uparrow$ & MT-B \\
\midrule
Default       & 0.30 & 0.20 & 0.50 & 0.315 & 0.615 & 0.660 & 7.29 \\
R1-dominant   & 0.60 & 0.10 & 0.30 & 0.362 & 0.562 & 0.602 & 7.05 \\
Equal         & 0.33 & 0.33 & 0.34 & 0.328 & 0.601 & 0.641 & 7.22 \\
R3-dominant   & 0.10 & 0.10 & 0.80 & 0.318 & 0.619 & 0.656 & 7.36 \\
R2-enhanced   & 0.20 & 0.40 & 0.40 & 0.341 & 0.588 & 0.628 & 7.15 \\
\bottomrule
\end{tabular}
\end{table}

Performance degrades monotonically as R1 weight increases: R3-dominant ($w_1 = 0.10$) nearly matches pure R3-only, while R1-dominant ($w_1 = 0.60$) is the worst configuration (NoF 0.362, WRR 0.562, MT-Bench 7.05). The optimal weight for R1 trends toward zero. WRR range across configurations is 5.7pp, indicating that the reward weights create a meaningful tradeoff, but the optimal operating point for every metric coincides with minimal R1 contribution. The R2-enhanced configuration ($w_2 = 0.40$) also underperforms, suggesting that the rank-based contrastive reward does not add value beyond what R3 captures.

\section{Discussion}

\paragraph{The direction is real but redundant.}
Our experiments confirm that a linear direction discriminating sycophantic from genuine responses exists in Qwen3-8B's mid-layer activations. The layer progression from chance (L0, AUROC 0.520) to peak (L19, AUROC 0.801) provides a meaningful window into where evaluative processing occurs. The interaction with evidence strength (weak: 0.773, strong: 0.831) reveals that the model's internal sycophancy encoding is weaker precisely when the task is most ambiguous. These are genuine findings about the model's representational structure.

However, this directional information does not translate into training utility. DC-GRPO offers no significant advantage over R3-only behavioral training. The continuous probe, which extracts the same information in a simpler form, achieves equivalent results. R1-R3 disagreement is too low (3.4\%) for the direction to serve a diagnostic role. The direction's content is largely captured by behavioral NLI scoring.

\paragraph{Why does the direction fail as a reward?}
Three factors contribute. First, at AUROC 0.801 with a wide CI (0.766--0.836), the signal-to-noise ratio is insufficient for reliable reward computation on individual samples. Second, direction drift during training (cosine similarity dropping to 0.66 in 3000 steps) means the reward signal is frequently stale. Third, R1 and R3 become highly correlated during training ($r > 0.9$ after step 4000), indicating that the model's strategy for satisfying R3 (behavioral correctness) automatically satisfies R1 (representational alignment) as a side effect.

\paragraph{Undifferentiated vs.\ targeted interventions.}
The one clear positive result is that targeted methods significantly outperform undifferentiated activation consistency (ACT-style). ACT penalizes all representational changes, including those corresponding to warranted belief revision, resulting in lower WRR (0.571 vs.\ 0.618--0.625). This confirms the double-bind hypothesis: effective anti-sycophancy training must be selective. But this selectivity can be achieved through behavioral NLI alone, without direct access to activation geometry.

\paragraph{Implications for RepE-as-reward.}
Our results suggest a boundary condition for using Representation Engineering directions as RL rewards: the approach requires both high discriminability ($\gg 0.80$) and stability under the policy changes induced by training itself. When discriminability is moderate and the direction drifts, the engineering overhead (calibration data construction, layer selection, recalibration scheduling, weight tuning) is not justified. These constraints may be specific to GRPO-style training, where large policy gradient updates rapidly reshape the activation landscape. Supervised fine-tuning with smaller update steps might preserve direction stability, though we leave this investigation to future work.

\paragraph{Cross-model structure.}
The finding that both Qwen3-8B and Llama-3-8B exhibit a sycophancy direction with similar layer-wise emergence patterns (peak at L19 and L18 respectively), but with nearly orthogonal specific directions (cosine 0.31), suggests that sycophancy encoding in mid-layers may be a general architectural property, while the specific geometry is model-family-specific. This limits the transferability of directional training methods but supports model-general interpretability studies.

\paragraph{Limitations.}
Our study is limited to 8B-parameter models and three random seeds, providing limited statistical power ($\alpha = 0.05$ tests with 3 seeds have low power for detecting small effects). The calibration data comprises 200 pairs, which may be insufficient for robust direction estimation. We evaluate on multi-turn dialogue sycophancy only; single-turn or domain-specific sycophancy may have different representational properties. Direction recalibration uses a fixed threshold strategy; adaptive methods might improve stability.

\section{Conclusion}

We investigated whether directional geometry in LLM activation spaces can serve as a training reward signal for reducing multi-turn sycophancy. While a sycophancy direction \dsyc{} exists with moderate discriminability and an informative layer-wise emergence pattern, it does not provide independent training value beyond behavioral NLI rewards when used within GRPO. Direction drift, limited signal-to-noise ratio at the critical weak-evidence condition, and high R1-R3 correlation during training collectively explain this negative result. A simpler continuous probe achieves equivalent outcomes without the engineering complexity of directional decomposition. These findings delineate practical boundary conditions for representation-level training rewards and redirect attention toward the conditions under which such signals might succeed: higher discriminability, training procedures that preserve representational stability, and settings where behavioral and representational signals genuinely diverge.

\bibliographystyle{plainnat}
\bibliography{references}

\end{document}
